\documentclass[10pt,a4paper]{amsart}
\usepackage{amsmath,amssymb,amsthm,graphicx}
\usepackage{verbatim}
\newenvironment{reference}{\comment}{\endcomment}
%\newenvironment{reference}{}{}
\newenvironment{slogan}{\comment}{\endcomment}
\newenvironment{history}{\comment}{\endcomment}

% For commutative diagrams we use Xy-pic
\usepackage[all]{xy}

% We use 2cell for 2-commutative diagrams.
\xyoption{2cell}
% Theorem environments.
%
\theoremstyle{plain}
\newtheorem{theorem}[subsection]{Theorem}
\newtheorem{proposition}[subsection]{Proposition}
\newtheorem{lemma}[subsection]{Lemma}

\theoremstyle{definition}
\newtheorem{definition}[subsection]{Definition}
\newtheorem{example}[subsection]{Example}
\newtheorem{exercise}[subsection]{Exercise}
\newtheorem{situation}[subsection]{Situation}

\theoremstyle{remark}
\newtheorem{remark}[subsection]{Remark}
\newtheorem{remarks}[subsection]{Remarks}

\numberwithin{equation}{subsection}

% Macros
%
\def\lim{\mathop{\mathrm{lim}}\nolimits}
\def\colim{\mathop{\mathrm{colim}}\nolimits}
\def\Spec{\mathop{\mathrm{Spec}}}
\def\Hom{\mathop{\mathrm{Hom}}\nolimits}
\def\Ext{\mathop{\mathrm{Ext}}\nolimits}
\def\SheafHom{\mathop{\mathcal{H}\!\mathit{om}}\nolimits}
\def\SheafExt{\mathop{\mathcal{E}\!\mathit{xt}}\nolimits}
\def\Sch{\mathit{Sch}}
\def\Mor{\mathop{\mathrm{Mor}}\nolimits}
\def\Ob{\mathop{\mathrm{Ob}}\nolimits}
\def\Sh{\mathop{\mathit{Sh}}\nolimits}
\def\NL{\mathop{N\!L}\nolimits}
\def\CH{\mathop{\mathrm{CH}}\nolimits}
\def\proetale{{pro\text{-}\acute{e}tale}}
\def\etale{{\acute{e}tale}}
\def\QCoh{\mathit{QCoh}}
\def\Ker{\mathop{\mathrm{Ker}}}
\def\Im{\mathop{\mathrm{Im}}}
\def\Coker{\mathop{\mathrm{Coker}}}
\def\Coim{\mathop{\mathrm{Coim}}}

% Boxtimes
%
\DeclareMathSymbol{\boxtimes}{\mathbin}{AMSa}{"02}

%
% Macros for moduli stacks/spaces
%
\def\QCohstack{\mathcal{QC}\!\mathit{oh}}
\def\Cohstack{\mathcal{C}\!\mathit{oh}}
\def\Spacesstack{\mathcal{S}\!\mathit{paces}}
\def\Quotfunctor{\mathrm{Quot}}
\def\Hilbfunctor{\mathrm{Hilb}}
\def\Curvesstack{\mathcal{C}\!\mathit{urves}}
\def\Polarizedstack{\mathcal{P}\!\mathit{olarized}}
\def\Complexesstack{\mathcal{C}\!\mathit{omplexes}}
% \Pic is the operator that assigns to X its picard group, usage \Pic(X)
% \Picardstack_{X/B} denotes the Picard stack of X over B
% \Picardfunctor_{X/B} denotes the Picard functor of X over B
\def\Pic{\mathop{\mathrm{Pic}}\nolimits}
\def\Picardstack{\mathcal{P}\!\mathit{ic}}
\def\Picardfunctor{\mathrm{Pic}}
\def\Deformationcategory{\mathcal{D}\!\mathit{ef}}
\usepackage{cleveref}
\crefname{originalreference}{original reference}{original references}
\Crefname{originalreference}{Original reference}{Original references}
\makeatletter
\newlabel{lemma-algebraic-extension-self-map}{{0BMD}{source}{}{sourcepointer-1}{}}
\newlabel{lemma-algebraic-extension-self-map@cref}{{[originalreference][0][]0BMD}{[1][0][]source}{}{}{}}
\makeatother
\crefname{originalsection}{original section}{original sections}
\Crefname{originalsection}{Original section}{Original sections}
\makeatletter
\newlabel{topology-example-automorphisms-of-a-set}{{0BMC}{source}{}{sourcepointer-2}{}}
\newlabel{topology-example-automorphisms-of-a-set@cref}{{[originalsection][0][]0BMC}{[1][0][]source}{}{}{}}
\makeatother
\crefname{originallemma}{original lemma}{original lemmas}
\Crefname{originallemma}{Original lemma}{Original lemmas}
\makeatletter
\newlabel{topology-lemma-profinite-group}{{0BR1}{source}{}{sourcepointer-3}{}}
\newlabel{topology-lemma-profinite-group@cref}{{[originallemma][0][]0BR1}{[1][0][]source}{}{}{}}
\makeatother
\crefname{originalsection}{original section}{original sections}
\Crefname{originalsection}{Original section}{Original sections}
\makeatletter
\newlabel{topology-section-topological-groups}{{0B1Y}{source}{}{sourcepointer-4}{}}
\newlabel{topology-section-topological-groups@cref}{{[originalsection][0][]0B1Y}{[1][0][]source}{}{}{}}
\makeatother
\title{The Galois group with its action topology is a profinite topological group}
\author{The Stacks Project Authors}
\date{stacks-project:fields.tex:local-snapshot-20261012; GNU FDL 1.2 or later}
\begin{document}
\maketitle
\noindent\textit{Editorial note.} All delivered mathematical text is the author's own, copied from the pinned Stacks Project chapter file fields.tex; no statement or proof step was rewritten, reordered or removed.\par\medskip
\noindent\textit{Editorial note.} The counted proof is the author's complete argument: The argument is nonstandard by design: the topology is defined before the inverse limit presentation, so profinite-ness is proved by exhibiting Gal(E/F) as a closed subset of the explicit profinite limit G = lim G(S).\par\medskip
\noindent\textit{Editorial note.} The label lemma-galois-profinite names the counted statement, and the proof environment immediately following it in the source is delivered verbatim as its proof.\par\medskip
\noindent\textit{Editorial note.} The counted statement is carried by the author's own lemma environment in the chapter source.\par\medskip
\noindent\textit{Allowed standard prerequisites (editorial).}\begin{itemize}
\item Topology, Example topology-example-automorphisms-of-a-set: the compact open topology on Map(E, E) is the universal topology making the action Map(E, E) x E -> E continuous (E discrete), and Aut(E) with the compact open topology is a topological group; cited as standard facts of the sibling Topology chapter
\item Topology, Lemma topology-lemma-profinite-group: a topological group whose underlying space is profinite is a profinite group
\item Same chapter: lemma-algebraic-extension-self-map (an F-algebra self-map of an algebraic extension is an automorphism)
\end{itemize}\medskip
\ifdefined\hypertarget\hypertarget{sourcepointer-1}{}\fi
\noindent\textit{Source reference.} Original reference 0BMD: sources/stacks-fields/pinned-original-source.tex, source lines 1024--1024. Prerequisite pointer invoked at 3159; it names a standard fact of the same chapter that the counted argument uses, and every step of the counted proof built on it is retained verbatim in the delivered body.\par\medskip
\ifdefined\hypertarget\hypertarget{sourcepointer-2}{}\fi
\noindent\textit{Source reference.} Original section 0BMC: sources/stacks-chow/tags.txt, source lines 14625--14625. Prerequisite pointer invoked at 3097; it names a standard fact of a sibling chapter that the counted argument uses, and every step of the counted proof built on it is retained verbatim in the delivered body.\par\medskip
\ifdefined\hypertarget\hypertarget{sourcepointer-3}{}\fi
\noindent\textit{Source reference.} Original lemma 0BR1: sources/stacks-chow/tags.txt, source lines 14754--14754. Prerequisite pointer invoked at 3122; it names a standard fact of a sibling chapter that the counted argument uses, and every step of the counted proof built on it is retained verbatim in the delivered body.\par\medskip
\ifdefined\hypertarget\hypertarget{sourcepointer-4}{}\fi
\noindent\textit{Source reference.} Original section 0B1Y: sources/stacks-chow/tags.txt, source lines 13911--13911. Prerequisite pointer invoked at 3113; it names a standard fact of a sibling chapter that the counted argument uses, and every step of the counted proof built on it is retained verbatim in the delivered body.\par\medskip
\begin{lemma}
\label{lemma-galois-profinite}
Let $E/F$ be a Galois extension. Endow $\text{Gal}(E/F)$ with the coarsest
topology such that
$$
\text{Gal}(E/F) \times E \longrightarrow E
$$
is continuous when $E$ is given the discrete topology. Then
\begin{enumerate}
\item for any topological space $X$ and map $X \to \text{Gal}(E/F)$
such that the action $X \times E \to E$ is continuous the induced map
$X \to \text{Gal}(E/F)$ is continuous,
\item this topology turns $\text{Gal}(E/F)$ into
a profinite topological group.
\end{enumerate}
\end{lemma}

\begin{proof}
Throughout this proof we think of $E$ as a discrete topological space.
Recall that the compact open topology on the set of self maps
$\text{Map}(E, E)$ is the universal topology such that the action
$\text{Map}(E, E) \times E \to E$ is continuous. See
Topology, Example \ref{topology-example-automorphisms-of-a-set}
for a precise statement. The topology of the lemma on
$\text{Gal}(E/F)$ is the induced topology coming from the
injective map $\text{Gal}(E/F) \to \text{Map}(E, E)$.
Hence the universal property (1) follows from the corresponding
universal property of the compact open topology.
Since the set of invertible self maps $\text{Aut}(E)$
endowed with the compact open topology forms a topological group, see
Topology, Example \ref{topology-example-automorphisms-of-a-set},
and since $\text{Gal}(E/F) = \text{Aut}(E/F) \to \text{Map}(E, E)$
factors through $\text{Aut}(E)$ we obtain a topological group.
In other words, we are using the injection
$$
\text{Gal}(E/F) \subset \text{Aut}(E)
$$
to endow $\text{Gal}(E/F)$ with the induced structure of a topological group
(see Topology, Section \ref{topology-section-topological-groups})
and by construction this is the coarsest structure of a topological
group such that the action $\text{Gal}(E/F) \times E \to E$ is continuous.

\medskip\noindent
To show that $\text{Gal}(E/F)$ is profinite we argue as follows
(our argument is necessarily nonstandard because we have defined
the topology before showing that the Galois group is an inverse
limit of finite groups).
By Topology, Lemma \ref{topology-lemma-profinite-group}
it suffices to show that the underlying
topological space of $\text{Gal}(E/F)$ is profinite.
For any subset $S \subset E$ consider the set
$$
G(S) = \{ f : S \to E \mid
\begin{matrix}
f(\alpha)\text{ is a root of the minimal polynomial}\\
\text{of }\alpha\text{ over }F\text{ for all }\alpha \in S
\end{matrix}
\}
$$
Since a polynomial has only a finite number of roots we see that
$G(S)$ is finite for all $S \subset E$ finite. If $S \subset S'$
then restriction gives a map $G(S') \to G(S)$. Also, observe
that if $\alpha \in S \cap F$ and $f \in G(S)$, then $f(\alpha) = \alpha$
because the minimal polynomial is linear in this case.
Consider the profinite topological space
$$
G = \lim_{S \subset E\text{ finite}} G(S)
$$
Consider the canonical map
$$
c : \text{Gal}(E/F) \longrightarrow G,\quad
\sigma \longmapsto (\sigma|_S : S \to E)_S
$$
This is injective and unwinding the definitions the
reader sees the topology on $\text{Gal}(E/F)$ as defined above
is the induced topology from $G$. An element $(f_S) \in G$ is in
the image of $c$ exactly if
(A) $f_S(\alpha) + f_S(\beta) = f_S(\alpha + \beta)$ and
(M) $f_S(\alpha)f_S(\beta) = f_S(\alpha\beta)$ whenever
this makes sense (i.e.,
$\alpha, \beta, \alpha + \beta, \alpha\beta \in S$).
Namely, this means
$\lim f_S : E \to E$ will be an $F$-algebra map
and hence an automorphism by
Lemma \ref{lemma-algebraic-extension-self-map}.
The conditions (A) and (M) for a given triple
$(S, \alpha, \beta)$ define a closed subset of
$G$ and hence $\text{Gal}(E/F)$ is homeomorphic
to a closed subset of a profinite space and therefore
profinite itself.
\end{proof}

\end{document}
